\section{Angles, Rotation, and the Unit Circle}
\noindent Angles provide a description of rotation. On the unit circle, this description determines a point whose coordinates define the sine and cosine functions.


\medskip
A rotation can be specified by its angle, independently of the size of the circle on which it acts. The relation between angle and position leads to trigonometric coordinates and to models of periodic motion.

\noindent\textbf{Learning objectives.} The principal objectives are to
describe signed rotation, convert between degrees and radians, calculate arc length
and angular speed, determine exact unit-circle coordinates, and relate uniform
circular motion to periodic coordinate functions.

\subsection{Angles measure turning}
An angle records a rotation from one ray to another. Counterclockwise rotations are positive, and clockwise rotations are negative. A rotation can exceed one full revolution. For example, $450^\circ$ and $90^\circ$ have the same terminal ray.

\begin{definition}Two angles are \emph{coterminal} if they have the same initial and terminal rays. They differ by an integer multiple of $360^\circ$, or $2\pi$ radians.\end{definition}
\begin{example}A positive coterminal angle for $-50^\circ$ is $-50^\circ+360^\circ=310^\circ$.\end{example}


\medskip
In standard position, the initial ray coincides with the positive horizontal axis.
Positive angles correspond to counterclockwise rotation and negative angles to
clockwise rotation. The signed angle specifies the accumulated rotation, whereas
the terminal ray specifies only the final direction. Thus $30^\circ$,
$390^\circ$, and $-330^\circ$ point along the same terminal ray, although
their rotations are different.

\begin{center}
\begin{tikzpicture}[scale=1.0]
\foreach \x/\a/\lab in {0/50/$50^\circ$,4/-50/$-50^\circ$}{
\begin{scope}[shift={(\x,0)}]
\draw[->,gray] (-1.2,0)--(1.4,0);
\draw[->,gray] (0,-1.25)--(0,1.25);
\draw[courseblue,thick] (0,0) circle (1);
\draw[very thick,coursegold,-{Latex}] (0,0)--(\a:1);
\draw[->,coursegreen,thick] (.43,0) arc[start angle=0,end angle=\a,radius=.43];
\node at (0,-1.55) {\lab};
\end{scope}}
\end{tikzpicture}
\end{center}

\noindent\textbf{Worked example: multiple turns.}
Find an angle between $0^\circ$ and $360^\circ$ coterminal with
$-770^\circ$. Adding $360^\circ$ twice gives $-50^\circ$; adding it
once more gives $310^\circ$. The resulting angle lies in $[0^\circ,360^\circ)$ and has the same terminal ray.

\begin{practice}Sketch $-120^\circ$, $240^\circ$, and $600^\circ$.
Which terminal rays coincide? Justify the result using coterminal angles.\end{practice}

\subsection{Radian measure}
For a circle of radius $r$ and an arc of length $s$ subtended by a central angle, the magnitude of the angle in radians is $s/r$. For an oriented arc, the signed arc length determines the signed angle. A full circumference is $2\pi r$, so one revolution measures $2\pi$ radians:
\[
\boxed{\theta=\frac{s}{r}},\qquad 180^\circ=\pi\text{ radians}.
\]
\begin{example}Convert $150^\circ$ to radians: $150(\pi/180)=5\pi/6$. Convert $7\pi/4$ radians to degrees: $(7\pi/4)(180/\pi)=315^\circ$.\end{example}
\begin{practice}Convert $225^\circ$ to radians and $-\pi/3$ radians to degrees, and sketch their terminal rays.\end{practice}

\begin{center}
\begin{tikzpicture}[scale=1.05]
\draw[thick,courseblue] (0,0) circle (1.7);
\draw[very thick,coursegreen] (1.7,0) arc[start angle=0,end angle=60,radius=1.7];
\draw[thick,coursegold] (0,0)--(1.7,0) node[midway,below] {$r$};
\draw[thick,coursegold] (0,0)--(60:1.7);
\draw[->,coursegold] (.5,0) arc[start angle=0,end angle=60,radius=.5];
\node at (.75,.34) {$\theta$};
\node[coursegreen] at (30:2.15) {$s=r\theta$};
\node[below] at (0,-2.1) {Arc length and central angle};
\end{tikzpicture}
\end{center}


\medskip
For a fixed central angle, the intercepted arc length is proportional to the radius. Consequently, the ratio $s/r$ does not depend on the size of the circle.

\begin{center}
\begin{tabular}{c|c|c}
Degrees & Radians & Fraction of one turn\\ \hline
$30^\circ$ & $\pi/6$ & $1/12$\\
$45^\circ$ & $\pi/4$ & $1/8$\\
$60^\circ$ & $\pi/3$ & $1/6$\\
$90^\circ$ & $\pi/2$ & $1/4$\\
$180^\circ$ & $\pi$ & $1/2$\\
$270^\circ$ & $3\pi/2$ & $3/4$\\
$360^\circ$ & $2\pi$ & $1$
\end{tabular}
\end{center}

The conversion $180^\circ=\pi$ radians determines every entry in the table. In particular, $60^\circ=\pi/3$ radians.

\begin{example}
Convert a negative angle: $-210^\circ\cdot\pi/180^\circ=-7\pi/6$.
Its terminal ray coincides with $5\pi/6$, because adding $2\pi$
to $-7\pi/6$ gives $5\pi/6$.
\end{example}

\subsection{Arc length and angular speed}
For signed radian measure $\theta$, the directed arc length is $s=r\theta$; the corresponding distance traveled is $r|\theta|$. Over a time interval of length $t$, the average angular velocity is $\omega=\theta/t$ and the average tangential speed is $v=r|\omega|$. For uniform rotation, these quantities are constant.
\begin{example}A wheel of radius $0.4$ m turns through $3\pi/2$ radians. A point on its rim travels $s=0.4(3\pi/2)=0.6\pi$ m.\end{example}
\begin{practice}A Ferris wheel of radius $12$ m makes one revolution every $40$ s. Find its angular speed and a passenger's tangential speed.\end{practice}


\medskip
\noindent\textbf{Worked example: a carousel.}
A carousel makes one revolution every $24$ seconds.
A child sits $3$ meters from the axis, while another sits $5$ meters
from the axis. Both riders have the same angular speed,
\[
\omega=\frac{2\pi}{24}=\frac{\pi}{12}\ \text{rad/s}.
\]
Their tangential speeds differ:
\[
v_1=3\frac{\pi}{12}=\frac{\pi}{4}\ \text{m/s},
\qquad
v_2=5\frac{\pi}{12}=\frac{5\pi}{12}\ \text{m/s}.
\]
Tangential speed depends on the distance from the axis, whereas angular speed does not.

The arc-length formula requires radian measure. For example, a quarter-turn is $\pi/2$ radians, so its arc length is $\pi r/2$.

\begin{practice}
A wheel of radius $0.6$ m turns through $135^\circ$ in $3$ seconds.
Find the arc length, average angular speed, and average tangential speed.
Explain the units of all three answers.
\end{practice}

\subsection{Coordinates on the unit circle}
Take a circle of radius $1$ centered at the origin. The point at angle $\theta$, measured from the positive $x$-axis, has coordinates
\[
\boxed{(x,y)=(\cos\theta,\sin\theta)}.
\]
Cosine is the horizontal coordinate and sine is the vertical coordinate. Because the point lies on $x^2+y^2=1$, we get $\cos^2\theta+\sin^2\theta=1$.
\begin{center}
\begin{tikzpicture}[scale=2]
\draw[->,gray](-1.2,0)--(1.25,0) node[right]{$x$};
\draw[->,gray](0,-1.2)--(0,1.25) node[above]{$y$};
\draw[courseblue,thick](0,0)circle(1);
\draw[orange!80!black,thick](0,0)--(40:1);
\draw[dashed](40:1)--({cos(40)},0);
\draw[->](.3,0)arc(0:40:.3);
\node at (.42,.15){$\theta$};
\node[above right]at(40:1){$(\cos\theta,\sin\theta)$};
\end{tikzpicture}
\end{center}


\medskip
Rotating the point $(1,0)$ through $\theta$ produces
$(\cos\theta,\sin\theta)$. Both coordinates are signed quantities;
their signs depend on the quadrant containing the terminal ray.

\begin{center}
\begin{tikzpicture}[scale=1.2]
\draw[->,gray] (-1.55,0)--(1.6,0) node[right]{$x$};
\draw[->,gray] (0,-1.5)--(0,1.55) node[above]{$y$};
\draw[courseblue,thick] (0,0) circle (1.25);
\draw[very thick,coursegold] (0,0)--(135:1.25);
\draw[dashed,coursegreen] (135:1.25)--({1.25*cos(135)},0);
\draw[dashed,courseblue] (135:1.25)--(0,{1.25*sin(135)});
\fill[coursegold] (135:1.25) circle (.04);
\node[above left] at (135:1.25) {$(-\frac{\sqrt2}{2},\frac{\sqrt2}{2})$};
\node at (.65,.62) {\scriptsize QI};
\node at (-.62,.62) {\scriptsize QII};
\node at (-.62,-.62) {\scriptsize QIII};
\node at (.62,-.62) {\scriptsize QIV};
\end{tikzpicture}
\end{center}

\noindent\textbf{Worked example: a point in Quadrant III.}
At $225^\circ=5\pi/4$, the reference angle is $\pi/4$.
The two coordinates have equal absolute value, but both are negative:
\[
(\cos(5\pi/4),\sin(5\pi/4))
=\left(-\frac{\sqrt2}{2},-\frac{\sqrt2}{2}\right).
\]
The reference triangle determines the coordinate magnitudes; the quadrant determines their signs.

\begin{practice}
Without a calculator, describe the point on the unit circle at
$150^\circ$ and at $330^\circ$. Which coordinates have the same
absolute value? Which signs differ?
\end{practice}

\subsection{Reference angles and periodicity}
Special angles arise from right triangles:
\[
\begin{array}{c|cc}
\theta&\cos\theta&\sin\theta\\\hline
0&1&0\\
\pi/6&\sqrt3/2&1/2\\
\pi/4&\sqrt2/2&\sqrt2/2\\
\pi/3&1/2&\sqrt3/2\\
\pi/2&0&1
\end{array}
\]
The reference angle supplies the absolute values in every quadrant; the quadrant gives the signs.
\begin{example}At $5\pi/6$, the reference angle is $\pi/6$ and the point is in quadrant II, so $\cos(5\pi/6)=-\sqrt3/2$ and $\sin(5\pi/6)=1/2$.\end{example}
After one full revolution the point returns to its previous position:
\[
\sin(\theta+2\pi)=\sin\theta,\qquad
\cos(\theta+2\pi)=\cos\theta.
\]
This geometric periodicity will become the foundation of sinusoidal graphs.


\subsection{Where exact coordinates come from}
The standard exact values follow from the geometry of two special right triangles.

An isosceles right triangle with legs $1$ has hypotenuse $\sqrt2$.
Dividing both legs by $\sqrt2$ yields the point
$(\sqrt2/2,\sqrt2/2)$ at $45^\circ$.
Likewise, a $30^\circ$--$60^\circ$--$90^\circ$ triangle has side
lengths in the ratio $1:\sqrt3:2$. Dividing by its hypotenuse
produces the coordinates $(\sqrt3/2,1/2)$ for $30^\circ$ and
$(1/2,\sqrt3/2)$ for $60^\circ$.

\begin{center}
\begin{tikzpicture}[scale=.95]
\begin{scope}
\draw[very thick,courseblue] (0,0)--(2,0)--(2,2)--cycle;
\draw (1.7,0) rectangle (2,.3);
\node[below] at (1,0) {$2$};
\node[right] at (2,1) {$2$};
\node[above left] at (1,1) {$2\sqrt2$};
\node at (.42,.18) {\scriptsize $45^\circ$};
\end{scope}
\begin{scope}[shift={(5,0)}]
\draw[very thick,coursegreen] (0,0)--({sqrt(3)},0)--({sqrt(3)},1)--cycle;
\draw ({sqrt(3)-.22},0) rectangle ({sqrt(3)},.22);
\node[below] at ({sqrt(3)/2},0) {$\sqrt3$};
\node[right] at ({sqrt(3)},.5) {$1$};
\node[above left] at ({sqrt(3)/2},.5) {$2$};
\node at (.42,.16) {\scriptsize $30^\circ$};
\end{scope}
\end{tikzpicture}
\end{center}

On the unit circle, the hypotenuse is normalized to length $1$, and the remaining side lengths give the absolute coordinate values.

\subsection{From the circle to a wave}
For uniform rotation with angular velocity $\omega$ and initial angle zero, let
\[
\theta(t)=\omega t,\qquad
x(t)=\cos(\omega t),\quad y(t)=\sin(\omega t).
\]
The point traces a circle in the coordinate plane; each coordinate, regarded as a function of time, oscillates periodically.

\begin{center}
\begin{tikzpicture}[x=.77cm,y=.9cm]
\draw[->,gray] (-.3,0)--(7.1,0) node[right]{$t$};
\draw[->,gray] (0,-1.4)--(0,1.5) node[above]{$y$};
\draw[coursegreen,very thick,domain=0:6.28,samples=110] plot(\x,{sin(\x r)});
\draw[courseblue,very thick,dashed,domain=0:6.28,samples=110] plot(\x,{cos(\x r)});
\draw[dotted] (1.5708,-1.2)--(1.5708,1.2);
\draw[dotted] (3.14159,-1.2)--(3.14159,1.2);
\draw[dotted] (4.71239,-1.2)--(4.71239,1.2);
\node[below] at (1.5708,0) {$\pi/2$};
\node[below] at (3.14159,0) {$\pi$};
\node[below] at (4.71239,0) {$3\pi/2$};
\node[below] at (6.28319,0) {$2\pi$};
\node[coursegreen,right] at (6.3,.5) {\scriptsize $\sin t$};
\node[courseblue,right] at (6.3,1.05) {\scriptsize $\cos t$};
\end{tikzpicture}
\end{center}
These graphs represent coordinate values as functions of time, rather than the circular trajectory in the $xy$-plane.

\subsection{Example: a Ferris-wheel model}
A Ferris wheel has radius $8$ meters, its center is $10$ meters above
ground, and a passenger starts at the rightmost point. Suppose it
rotates counterclockwise once every $40$ seconds. Its angular speed is
$\omega=2\pi/40=\pi/20$ radians per second, and the passenger's
position relative to the ground is
\[
\boxed{
x(t)=8\cos\left(\frac{\pi t}{20}\right),\qquad
y(t)=10+8\sin\left(\frac{\pi t}{20}\right).
}
\]
At $t=0$ the passenger is at $(8,10)$. At $t=10$ seconds the
wheel has turned through $\pi/2$, so the passenger is at $(0,18)$.
At $t=20$ seconds the passenger is at $(-8,10)$. The coordinate equations determine the passenger's position at any time.

The radius determines the amplitude of both coordinate functions. The center height shifts the vertical coordinate, and the period determines the angular velocity.

\begin{practice}
A different wheel has radius $6$ meters, center height $7$ meters,
and period $30$ seconds. Starting from the rightmost point,
write its coordinate model and find the passenger's height after
$7.5$ seconds.
\end{practice}

\subsection{Exercises}
\begin{enumerate}
\item Convert $-135^\circ$ and $300^\circ$ to radians.
\item Convert $11\pi/6$ and $-5\pi/4$ to degrees.
\item Find an angle in $[0,2\pi)$ coterminal with $19\pi/6$.
\item A circle of radius $8$ cm subtends an angle of $5\pi/12$. Find the arc length.
\item A turntable makes $33\frac13$ revolutions per minute. Find angular speed in radians per second.
\item Explain why sine and cosine must lie between $-1$ and $1$.
\item Find unit-circle coordinates at $2\pi/3$, $5\pi/4$, and $11\pi/6$.
\item A Ferris wheel has radius $10$ m and center $(0,12)$. Write a rider's coordinates after turning through $\theta$ from the rightmost point.

\item A point begins at $(1,0)$ and turns through $-7\pi/6$ radians. Determine a coterminal angle in $[0,2\pi)$ and give the exact point.
\item A wheel completes $9$ revolutions in $2$ minutes. Determine its angular speed in both radians per second and degrees per second.
\item A rider on a circle of radius $4$ m moves through $75^\circ$. Find the arc length exactly in terms of $\pi$ and approximately to two decimals.
\item Describe the coordinates at $3\pi/2$ and explain why a terminal point on an axis does not belong to any open quadrant.
\item For the Ferris-wheel model in the chapter, explain which physical quantity would change if the coefficient $8$ in both coordinate equations became $12$.
\end{enumerate}
\noindent The unit-circle coordinates also determine tangent, cotangent, secant, and cosecant, which are developed in the following chapter.
\clearpage
